\section{THEORETICAL FRAMEWORK}
\label{sec:theory}

\subsection{Core Theoretical Constructs}

\subsubsection{The Transformational Container}

The concept of the transformational container originates in the psychoanalytic work of \textcite{winnicott1960theory} and was adapted to role-playing game theory by \textcite{bowman2021magic}. It describes a bounded space, building upon \textcite{huizinga1958homo}'s ``magic circle'' but additionally requiring consent, safety structures, and community support oriented toward lasting personal change \parencite{bowman2021magic, baird2022transformative, bowman2025transformative}. When these conditions hold, the container provides participants enough psychological safety to risk vulnerability; the resulting discomfort becomes the material for growth. I argue that debate, in its existing forms, already constitutes a partially-realized transformational container. The format's explicit rules and conventions (speaking times, turn order, permissible modes of engagement) establish boundaries that demarcate the debate encounter from ordinary social interaction \parencite{bowman2025transformative}. Within these boundaries, the competitive frame provides what \textcite{deterding2017alibis} terms \textit{alibi}: because side allocation is typically random and the activity is understood as a game, debaters can explore unfamiliar arguments and adopt perspectives they might not otherwise endorse with reduced social risk. The judge functions as the containing authority analogous to the facilitator in transformative game design \parencite{bowman2025transformative}: a figure who maintains the rules, calibrates the encounter through feedback, and structures the post-round evaluation.

% Crucially, the container is a \textit{group} space, and group dynamics shape its transformative capacity. \textcite{yalom2020theory} identifies therapeutic factors operating in group settings---universality (discovering others share one's struggles), group cohesion, altruism, and interpersonal learning---all of which are present in debate communities. The shared experience of intellectual vulnerability, of being publicly challenged and occasionally defeated, creates bonds that support continued risk-taking. When these group dynamics function well, the container is strengthened; when status hierarchies or competitive antagonism dominate, the group forces that sustain the container erode.

% \subsubsection{Identity and the Debater-Character Relationship}

% If the container provides the bounded space within which transformation occurs, identity theory specifies \textit{what} is being transformed. \textcite{bowman2010functions} proposes a model of identity in role-playing as fundamentally fragmented: rather than possessing a single, unified self, players bring multiple identity facets---values, beliefs, competencies, emotional dispositions---any of which may be activated, magnified, or recombined through character play. \textcite{bowman2025transformative} develop this further through the concept of \textit{mosaic identity}, where the self is understood as a composite of overlapping fragments that are differentially activated across contexts. Role-playing does not create identity from nothing; it selects, amplifies, and reorganizes existing facets of the player's self.

% The debater-character relationship exemplifies this fragmented self model. As established in the background chapter, debate characters represent fragmented selves whose defining feature is epistemic commitment to a particular position. The debater locates within themselves beliefs, knowledge, and rhetorical dispositions that can support the assigned side, magnifying these into the character's core identity while suppressing contradictory facets. This is not pretense but genuine identity work: the cognitive demands of competitive performance require authentic inhabitation of the epistemic stance, producing what \textcite{bowman2010functions} describes as developing a ``theory of mind'' outside one's own. \textcite{lankoski2012embodied} explain this through grounded cognition theory: performing an epistemic stance neurologically activates the same processes as holding that stance, meaning that the debater-character's identity work is cognitively genuine rather than merely theatrical.

% This identity framework is essential for understanding the bleed mechanisms that follow. Memetic bleed operates precisely through identity facets: arguments rehearsed and positions inhabited do not simply vanish when the round ends, because they were constructed from genuine fragments of the debater's self. The facets activated during play---the capacity to see merit in an opposing view, the rhetorical skill of making a case one does not initially endorse---persist and integrate into the debater's broader identity mosaic. Emancipatory bleed, similarly, works through identity: for marginalized students, the experience of being taken seriously as intellectuals activates and strengthens identity facets---confidence, intellectual authority, public voice---that may have been suppressed in other contexts.

\subsubsection{Bleed as Transformative Mechanism}

Following \textcite{hugaas2024bleed}, bleed describes the bidirectional spillover of psychological contents between player and character (in debate, between the debater and their assigned position). \textcite{leonard2018bleed}, drawing on \textcite{lankoski2012embodied}'s grounded cognition framework, provide neuroscientific grounding: embodied cognition means inhabiting a role activates processes that do not fully distinguish fictional from real experience, and the self-regulation required to compartmentalize in-character reactions depletes under sustained performance. Applied to debate, this suggests that extended rounds of inhabiting an assigned position gradually erode the cognitive separation between debater and debater-character, helping explain debate's lasting attitudinal effects. In debate, I identify 3 types of bleed at play.

Memetic bleed, the transfer of ideas and perspectives, explains how debaters develop what I call epistemic humility: skepticism towards one's own opinions and openness to the beliefs of others. The opinions of the role played in a debate bleed directly toward the debater, becoming arguments that can be summoned again outside the round. More importantly, after years of arguing positions contrary to their beliefs, debaters should bleed into accepting that reasonable people hold opposing views and that argumentative power does not equate with being right. This mechanism provides a theoretical account of what forensics scholars have documented empirically: Colbert's ``intellectual disequilibrium'' \parencite{colbert1995enhancing} describes the trigger that opens the channel for memetic bleed, while the ``democratic ethos'' that \textcite{hicks1998public} identifies as debate's civic outcome is the intended result of that bleed integrating into the debater's empowered attitude. This should render it impossible for value convergence, the process of members of a community becoming intolerant to other opinions, to emerge. Later on I hypothesize why I observed it doing so.

This bleed of a humble attitude connects to \textcite{bogost2007persuasive}'s concept of procedural rhetoric. Bogost argues that every game, through its systems and rules, reflects a world-view that players may in time be persuaded of. The rules of debate inherently postulate that anyone can support anything and that the dialectic process reaches not truth but a showcase of who is more capable in argument (the same motion can be decided either way in different rounds). Thus experienced debaters ideally learn to acknowledge how there may be much truth in the points made also by the less skilled in argument.

Procedural bleed, the transfer of physical abilities and habits, explains skill development. The rhetorical habits, speaking patterns, and argumentative reflexes developed through competition bleed into debaters' ordinary communication, producing the documented improvements in public speaking and critical thinking. \textcite{cirlin1997sociological}'s account of how debaters develop persuasive skill through audience adaptation describes precisely this form of procedural bleed and its transfer beyond the activity.

Emancipatory bleed \parencite{kemper2017battle,kemper2020wyrding} explains debate's particular power for students from marginalized backgrounds. For students where eloquent public speaking was not modeled, debate provides a transformational container \parencite{bowman2021magic} to discover and develop their voice. Being taken seriously as intellectuals, having ideas weighed on merits rather than dismissed based on identity, produces lasting changes in self-concept and capability. \textcite{snider2003gamemaster}'s documentation of debate empowering previously ``silenced student populations'' through Urban Debate Leagues provides empirical evidence for this form of bleed in competitive debate contexts.

%\subsubsection{Debrief and Integration Practices}

%\textcite{bowman2025transformative} present a taxonomy of debrief practices essential to transformative design, distinguishing several dimensions: structured versus unstructured debriefs, emotional versus intellectual versus educational processing, and practices of de-roling (explicitly stepping out of character), narrativization (constructing meaning from the experience), and post-game integration (connecting in-game experiences to out-of-game life). The debrief is not merely a postscript to play but the primary site where raw experience is converted into lasting transformation \parencite{crookall2014engaging,daniau2016transformative}.

%Traditional debate's built-in debrief---the oral critique---is revealing when analyzed through this taxonomy. It is structured but narrowly so: the judge explains their decision, identifies strengths and weaknesses, and offers advice for improvement. This processing is almost exclusively intellectual and evaluative. The judge assesses what worked and what did not as argumentation, with no attention to emotional processing, no de-roling (debaters are expected to simply move on to the next round), and minimal integration connecting the argumentative experience to personal development. The absence of emotional processing is particularly significant given that debate regularly requires inhabiting positions contrary to one's beliefs, engaging with emotionally charged topics, and experiencing public evaluation. \textcite{montola2010positive} argues that negative emotions in role-playing---discomfort, vulnerability, frustration---can be transformative rather than merely distressing, but only when the container provides adequate processing space. Traditional debate formats offer no such space: the oral critique processes the \textit{game performance} but not the \textit{human experience}.

%The RPG-specific debrief literature reinforces this concern: \textcite{atwater2016debrief}'s foundational review characterizes debriefs as ``non-therapeutic transitional tools for moving between perspectives'' and cautions that debriefing alone does not make activities safer---safety must be designed into the format itself. This analysis reveals debrief design as a major untapped lever for enhancing debate's transformative potential. Formats might incorporate emotional check-ins, explicit de-roling moments separating the debater from the position argued, or structured reflection connecting argumentative experiences to the debater's developing worldview.

%\subsubsection{Safety Frameworks}

%\textcite{bowman2025transformative} present a comprehensive safety model organized across three temporal phases and a community-level dimension. \textit{Before-game} safety includes briefing participants on content and expectations, establishing consent frameworks, and calibrating the intensity of play. \textit{During-game} safety encompasses real-time safety mechanics enabling participants to signal discomfort, escalate or de-escalate intensity, and exit if needed, along with facilitator monitoring of participant well-being. \textit{After-game} safety involves structured debrief, emotional check-ins, and aftercare practices supporting integration. The \textit{community-level} dimension addresses systemic safety: codes of conduct, reporting mechanisms, power dynamics awareness, and community culture. Complementing this temporal model, \textcite{bowman2025safety} propose a participant-experience framework of ``Zones of Safety, Challenge, and Risk'': the \textit{comfort zone}, where participants engage with familiar material; \textit{growth edges}, where productive discomfort enables learning; and \textit{hard limits}, where content risks genuine harm. Applied to debate, the comfort zone encompasses familiar topics and preferred argumentation styles; growth edges correspond to challenging motions and arguing unfamiliar sides; and hard limits involve topics touching personal trauma or identity threats.

%Analyzed through this model, competitive debate's safety provisions are unevenly developed. Before-game safety is minimal: debaters receive motions with little content warning, and there is no calibration of emotional intensity. During-game safety relies almost entirely on the competitive frame itself---the alibi that ``it's just a debate'' is expected to provide sufficient protection, but this alibi erodes precisely when topics become personally relevant or when competitive stakes are high. After-game safety, as discussed in the debrief section, processes game performance but not emotional experience. At the community level, while organizations like ARDOR maintain ethical guidelines and equity officers, the competitive structure itself produces hierarchies and pressures that can undermine psychological safety.

%This analysis is significant because the thesis identifies competitive pressure and hierarchical dynamics as central problems---and these are fundamentally safety concerns. The methods chapter already commits to safety practices including pre-session briefing, post-session debrief, equity officer access, and the ability to exit without explanation. The safety framework from \textcite{bowman2025transformative} provides the theoretical grounding for these design choices and suggests additional interventions: content calibration before rounds, in-game safety mechanics adapted from larp practice, and structured emotional processing after rounds.

%\subsubsection{Designable Surfaces}

%A core design principle articulated in \textcite{bowman2025transformative} holds that in transformative game design, \textit{everything is a designable surface}: every rule, role, ritual, spatial arrangement, temporal structure, and social norm shapes participants' experience and therefore their potential for transformation. This principle reframes the designer's task from creating a single ``game'' to intentionally shaping a constellation of interconnected surfaces, each of which can be independently tuned to support or undermine transformative goals.

%Applied to competitive debate, the designable surfaces include: the motion or topic (what participants argue about, and how it is framed); side assignment (whether debaters choose or are assigned positions, and how this structures their relationship to the content); speech structure and timing (the sequence, duration, and format of speeches, which shape cognitive load and interaction patterns); the judge's role and evaluative criteria (who the judge role-plays as, and what standards they apply); feedback and debrief structure (how post-round processing is organized, as discussed above); competitive structure (how winners are determined, how status is distributed, what incentives the tournament creates); and community norms (the informal expectations, traditions, and social dynamics surrounding the formal game).

%Current debate formats have evolved these surfaces primarily through tradition, convention, and incremental competitive optimization rather than through intentional transformative design. The World Schools format's speech timings, for instance, derive from practical compromise rather than from analysis of optimal cognitive load for perspective-taking. Judging criteria emphasize competitive fairness---a worthy goal---but do not explicitly consider how evaluative standards shape the bleed patterns discussed above. This principle provides the theoretical warrant for the thesis's design project: if every surface shapes transformation, then every surface is an opportunity for intentional intervention.

\subsubsection{Zone of Proximal Development and Scaffolded Learning}

Vygotsky argues that higher cognitive functions develop through the scaffolding, by a more capable other, of challenges. When learners operate in the zone of proximal development where challenge is achievable only with guidance, they internalize new competences. \textcite{vygotsky1978mind}'s framework explains how debate creates naturally occurring scaffolding: power-pairing keeps challenges within this zone, while judges provide the guided instruction that feeds internalization. The debate cycle naturally embodies \textcite{kolb1984experiential} experiential learning model, with the debrief, which \textcite{crookall2014engaging} and \textcite{daniau2016transformative} identify as essential to transformation, built into debate through the oral critique. This scaffolding explains why the ``intellectual disequilibrium'' that \textcite{colbert1995enhancing} identifies as debate's cognitive engine produces growth rather than frustration: coaching, graduated competition, and structured feedback keep the cognitive challenge within the zone where productive accommodation occurs.

\subsection{Diagnosing Forensics Pedagogy Critiques through Transformative Game Design Mechanisms}

The critiques identified within forensics pedagogy, echoing the issues identified during the problem statement, can be productively reframed through the theoretical mechanisms of transformative game design, illuminating not only why these problems occur but also suggesting design-based interventions.

\subsubsection{Value Convergence as Corrupted Memetic Bleed}

What I coined as value convergence can be seen as a failure in memetic bleed: the format-level outcome it suppresses is \textit{value pluralism}, and the participant-level outcome is \textit{epistemic humility} (the same transformation seen from two sides). Whereas forensics literature locates the cause in a ``closed system'' of insulated tournament practice, the role-playing-game framing allows me to isolate a clear mechanism within that system: standard formats ask every judge to role-play the same archetype, the average unbiased person. This role is vaguely specified, conflating ``unbiased'' with ``lacking values'' and thus procedurally teaching that the unspoken values that debaters try to converge to in an attempt to make the game more objective and competitively unfrastrating are an acceptable baseline of society and thus correct.

Following \textcite{bowman2025transformative}, the container's effectiveness requires psychological safety for all participants. When judges face social consequences for heterodox opinions the alibi protecting honest judgment erodes and what bleeds into debaters is strategic equilibrium rather than appreciation for intellectual diversity. Debaters, optimizing for competitive success, do ``argument borrowing'', adapting strategically: arguments that succeed in elimination rounds or from high-status debaters become templates, and approaches aligned with perceived values propagate. These two effects in time lead community members to see said values as the only acceptable ones, integrating them in their identity. 

Design interventions might therefore (a)~diversify audiences debaters must persuade, (b)~strengthen judge alibi through structural features, and (c)~explicitly distinguish in debriefing between ``what won this debate'' and ``what constitutes good reasoning.''

\subsubsection{Information Density and Misaligned Procedural Bleed}

The critique of information-dense speaking represents misalignment in procedural bleed. When formats reward rapid delivery optimized for expert judges, debaters develop procedural habits adapted to highly specific audiences. The skills that bleed out are real but narrowly optimized, representing failure of transfer. Rubric-perception asymmetry creates the condition for misaligned bleed; the judge-as-expert role converts it into a bleed pattern; community governance propagates that pattern across a circuit. I will refer to the positive transformation along this axis as \textit{transferable persuasion}: rhetorical habits that survive the move from expert audiences to ordinary ones.

Design interventions might include (a)~evaluation criteria rewarding accessibility alongside substance, or (b)~judging paradigms simulating general audiences. (c)~By constructing vivid characters for judges to role-play, the bleed-in from character to judge can itself become memetic, the judge taking on the character's understanding while maintaining expert-level rule comprehension.

\subsubsection{Competitive Hierarchy as Compromised Container}

Hierarchy is downstream of a deeper mechanism. As foreshadowed in §3.1.1, sustained competitive pressure erodes the alibi enabling transformation: when tournament success conflates with status and self-worth, debate can no longer be treated as ``just a game,'' and stable hierarchies form along the axis of who has won. I will refer to the positive outcome along this axis as \textit{container safety}: a play environment whose stakes do not collapse into personal worth. \textcite{bowman2013social} documents social conflict within role-playing communities arising from power differentials between facilitators and players, competitive play dynamics, creative agenda disputes, and unprocessed bleed. These patterns parallel competitive debate: power differentials map onto hierarchies between experienced and novice debaters and competitive play produces status systems based on tournament success.

Emancipatory bleed requires a container where being taken seriously is possible; rigid hierarchies block this for those positioned as subordinate. The zone of proximal development requires calibrated challenges; competitive anxiety may prevent productive struggle or reflective observation.

Design interventions might (a)~distribute status more broadly, (b)~emphasize growth over victory, (c)~strengthen alibi by making explicit the distinction between game performance and personal worth, and (d)~incorporate safety practices adapted from larp design.

%\subsubsection{Implications for Design}

%This reframing generates specific predictions. Formats designed to maximize transformation should: (1) structure memetic bleed toward epistemic humility through debriefing practices connecting argumentative experiences to personal development; (2) align procedural bleed with transferable communication skills by strengthening judge alibi to role-play as specific general audience members; and (3) maintain container safety by preventing hierarchical success from conflating with personal worth.
